\scanpage{7}
On trouve en particulier

\uline{Proposition} Les formes bilinéaires \uline{symétriques} $\varphi$ sur $E$ définissent des structures d'\uline{Algèbre de Lie} sur $E'$.

\uline{Exemple} Prenons $E$ muni de la base (8), \deleted{\unsure{et envi-}}\inserted{d'où} identification
\[
 E\simeq\underline{\mathrm O}_S^3 \tag*{(20)}
\]
\unsure{écrivons|décrivons|choisissons|dirons} sur $E$ la forme quadratique
\[
 q(x,y,z)=xy+z^2 \tag*{(21)}
\]
correspondant à la forme bilinéaire associée
\[
 \varphi((x,y,z)(x',y',z'))=(xy'+yx')+2zz' \tag*{(22)}
\]
donc
\[
\left\{\begin{aligned}
 q(e_1)=q(e_2)&=0,\quad q(e_3)=1\\
 \varphi(e_1,e_2)&=1,\quad \varphi(e_2,e_3)=\varphi(e_3,e_1)=0\\
 \varphi(e_1,e_1)&=\varphi(e_2,e_2)=0\qquad\varphi(e_3,e_3)=2
\end{aligned}\right. \tag*{(23)}
\]
\margin{\uline{NB} $\varphi(\xi,\xi)=2q(\xi)$}
(ou encore
\[
\left\{\begin{aligned}
 a=b&=0,\quad c=1\\
 p=p'=q=q'&=0\qquad r=r'=1
\end{aligned}\right.,\quad \text{matrice})\ ,\varphi=\begin{pmatrix}0&1&0\\1&0&0\\0&0&2\end{pmatrix} \tag*{(23 bis)}
\]
Les formules (12) deviennent alors \deleted{\unsure{ainsi}}
\[
\left\{\begin{aligned}
 [e'_2\unsure{\wedge|,}e'_3]&=e'_2\\
 [e_3,e'_1]&=e'_1\\
 [e'_1,e'_2]&=2e'_3
\end{aligned}\right. \tag*{(24)}
\]
\margin{NB $\begin{aligned}f(e_1)&=e'_2\\ f(e_2)&=e'_1\\ f(e_3)&=2e'_3\end{aligned}$}
ou encore, posant
\[
 X=e'_1,\quad Y=e'_2,\quad H=e'_3 \tag*{(25)}
\]
on reconnait les formules familières
\[
 [H,X]=X,\quad[H,Y]=-Y,\qquad[X,Y]=2H \tag*{(26)}
\]
de l'algèbre de Lie \deleted{\unsure{du|de}} ``universelle'' (sur les entiers) de $\mathrm{GP}(1)$.

% archive sheet number: 6
